\documentclass[10pt,a4paper]{amsart}
\usepackage[margin=19mm]{geometry}
\usepackage{amsmath,amssymb,mathtools}
\usepackage[scr=rsfs,cal=euler]{mathalfa}
\usepackage{xfrac}
\usepackage[unicode,hidelinks,bookmarks=false]{hyperref}
\newcommand{\ie}{i.e.\ }
\newcommand{\cf}{cf.\ }
\newcommand{\resp}{respectively\ }
\newcommand{\Subsection}{\subsection}
\providecommand{\nobreakdash}{\nobreak\hbox{-}}
\setlength{\emergencystretch}{2em}
\multlinegap0pt
\usepackage{amssymb}
\usepackage{xcolor}
\newtheorem{cor}{Corollary}
\newtheorem{thm}{Theorem}
\newtheorem{prop}{Proposition}
\title{Segmentation of monotone data by Kobayashi-Warren-Carter type total variation energies}
\author{Yoshikazu Giga, Ayato Kubo, Hirotoshi Kuroda, Koya Sakakibara}
\date{Source: arXiv:2603.28078v2 (CC BY 4.0)}
\begin{document}
\maketitle

\noindent\emph{Source and licence.} Segmentation of monotone data by Kobayashi-Warren-Carter type total variation energies. Version arXiv:2603.28078v2 (CC BY 4.0), distributed by arXiv under the Creative Commons Attribution 4.0 International licence, http://creativecommons.org/licenses/by/4.0/, as recorded in the arXiv licence metadata for this exact version and shown on the abstract page https://arxiv.org/abs/2603.28078v2. Full licence notice: \texttt{licenses/arxiv-2603.28078-CC-BY-4.0.txt}.\par\smallskip

\noindent\textit{Editorial scope.} Counted unit: the article's thm TCJ (source0.tex lines 523--545). The article's complete original proof of it is delivered with it (source0.tex lines 546--678, one \texttt{proof} environment of 133 lines). No mathematics is written, repaired or re-derived: the statement and the proof are the article's own lines, in the article's own notation, and no retained line is substituted or re-flowed.  Retained uncounted as the source-local context the proof needs: lines 417--421; lines 680--690.  One declared adaptation: exactly 1 ancillary strings inside the retained spans are omitted (image-inclusion commands and, where the source repeats a label, the repeated label definitions) and no other line differs from the source; the figure environments, with their placement, captions and labels, are kept, and the package records the omitted strings in fidelity receipts bound to the affected bodies.  No retained line contains a citation, so the wrapper carries no bibliography and no bibliography entry.  The retained lines contain the article's own diagram code, which the wrapper typesets with the standard tikz packages; no image file is delivered and the package declares no asset dependency.  Editorial formatting only, in addition: an AMS \texttt{amsart} preamble that restates the article's own declarations and macros used by the retained lines, namely the environments \texttt{cor}, \texttt{thm}, \texttt{prop} on separate counters, together with the standard packages those lines need.  The retained environments print in this document as Corollary 1, Theorem 1, Proposition 1 in order of appearance, whereas the article prints the same items with the numbering of its own full document.  The complete native source of the article as distributed in the arXiv e-print for this exact version is bundled unchanged as \texttt{sources/197408/source0.tex}.
\begin{cor} \label{CApMon}
Assume that $K$ satisfies (K1), (K3) and that $\Omega=(a,b)$.
 Assume that $u$ is non-decreasing (resp.\ non-increasing) and bounded in $(a,b)$.
 Then there is a sequence of non-decreasing (non-increasing) piecewise constant functions $\{u_m\}$ (with finitely many jumps) such that $u_m\to u$ in $L^p(\Omega)$ for any $p\ge1$ and $TV_K(u_m)\to TV_K(u)$ as $m\to\infty$.
\end{cor}

\begin{prop} \label{PInd}
Assume that $K$ satisfies (K2).
 Let $M>0$.
 If $\rho_i\ge0$ ($0\le i \le k$) satisfies $\sum_{i=0}^k\rho_i\le M$, then
\begin{equation} \label{EInd}
    \sum_{i=0}^k K(\rho_i)
    \ge K \left(\sum_{i=0}^k \rho_i \right)
    + C_M \sum_{0\le i<\ell\le k} \rho_i \rho_\ell
\end{equation}
for $k\ge1$.
\end{prop}

\begin{thm} \label{TCJ}
Assume the same hypotheses of Corollary~\ref{CApMon} concerning $K$, $\Omega$, $u$ and $u_m$.
 Assume further (K2).
 Then for each jump interval $I_{\lambda_j}(\rho_j)$, there exists a subsequence $\{i_m^j\}_{m=1}^\infty\subset\mathbb{N}$ such that $\left\{I_{\lambda_{i_m^j}^m}(\rho_{i_m^j}^m)\right\}_{m=1}^\infty$ converges to $I_{\lambda_j}(\rho_j)$ as $m\to\infty$.
 In other words,
\[
    \lambda_{i_m^j}^m \to \lambda_j, \quad
    \rho_{i_m^j}^m \to \rho_j
    \quad\text{as}\quad m \to \infty.
\]
Moreover, if a subsequence $\{i_m\}$ is taken such that
\[
    \lambda_{i_m}^m \to \lambda, \quad
    \rho_{i_m}^m \to \rho
    \quad\text{as}\quad m \to \infty
\]
with some $\lambda$ and $\rho$ and $\lambda\neq\lambda_j$ for any $j$, then $\rho=0$.
 The convergence $\rho_{i_m}^m\to0$ is uniform, i.e., 
\[
    \lim_{m\to\infty} \max \left\{\; \rho_\ell^m \bigm|
    \ell\neq i_m^j\ \text{for all}\ j\; \right\} = 0.
\] 
\end{thm}

\begin{proof}
Since $u_m\to u$ in $L^1$, $u_m^{-1}\to u^{-1}$ in $L^1(G)$ as $m\to\infty$ for any closed interval $G \subset\left(u(a+0),u(b-0)\right)$.
 By definition, $u^{-1}=z_j$ in $I_{\lambda_j}(\rho_j)$ and $u^{-1}(p)<z_j$ for $p<\lambda_j$ and $u^{-1}(p)>z_j$ for $p>\lambda_j+\rho_j$.
 The $L^1$-convergence implies that there is a subsequence $\{i_m^j\}$, $\{\bar{i}_m^j\}$ with $i_m^j\le \bar{i}_m^j$ such that
\[
    p_m = u_m(z_{i_m^j}^m) \to \lambda_j \quad
    q_m := u_m(z_{\bar{i}_m^j}^m + 0) \to \lambda_j + \rho_j
\]
as $m\to\infty$.

We consider $TV_K$ near $z_j$.
 For $\varepsilon>0$, we take $\delta>0$ small such that
\begin{equation} \label{EATVK}
    TV_K(u,Z_j^\delta) \le K(\rho_j) + \varepsilon, \quad
    Z_j^\delta = (z_j-\delta, z_j+\delta).
\end{equation}
% 原稿 2026/1/21、6/16
By definition, on the interval
\[
    (p'_m, q'_m]
    = \left( u_m(z_{i'_m}^m), u_m(z_{i'_m}^m + 0) \right],
\]
$u_m^{-1}$ equals a constant $z_{i'_m}^m$.
 Assume that $i_m^j\le i'_m\le \bar{i}_m^j$.
 Since $u_m^{-1}$ converges to $u^{-1}$ in $L^1$,
 we may assume that $p'_m$ and $q'_m$ converge to some $p$, $q$ as $m\to\infty$ with
\[
    \lambda_j \le p, \quad
    q \le \lambda_j + \rho_j;
\]
see Figure~\ref{FAp}.
\begin{figure}[tb]
\centering 

\caption{the graph of $u^{-1}$ and $u_m^{-1}$\label{FAp}}
\end{figure}

\noindent
\emph{Case} 1 ($p<q$).
 We shall prove that $\lambda_j=p$, $\rho_j=q-p$.
 We set $q_m-p_m=:\rho^m$, $q'_m-p'_m=:\rho'_m$.
 By definition of $Z_j^\delta$, for sufficiently large $m$ we see that
\[
    TV_K(u_m,Z_j^\delta) \ge K(\rho'_m)
    + K(\rho^m -\rho'_m).
\]
Since $TV_K(u_m)\to TV_K(u)$ as $m\to\infty$ implies
\[
    \lim_{m\to\infty} TV_K (u_m,Z_j^\delta)
    = TV_K (u,Z_j^\delta),
\]
by (K1), sending $m\to\infty$ yields
\[
    TV_K(u,Z_j^\delta)
    = \varliminf_{m\to\infty} TV_K(u_m,Z_j^\delta)
    \ge K((\rho_*) + K(\rho_j-\rho_*)
\]
where $\rho_*=\lim_{m\to\infty}\rho'_m$.
 By our assumption $\rho_*=q-p>0$ and (K2) we see that
\[
    K(\rho_*) + K(\rho_j-\rho_*)
    \ge K(\rho_j) + C_M \rho_*(\rho_j-\rho_*)
\]
where $M=\rho_j$.
 Thus,
\[
    TV_K(u,Z_j^\delta) \ge K(\rho_j)
    + C_M \rho_*(\rho_j-\rho_*).
\]
By \eqref{EATVK}, we conclude that $\rho_j-\rho_*=0$.
 Thus $\lambda_j=p$, $\rho_j=q-p$.

\noindent
\emph{Case} 2 ($p=q$).
 Assume that for all subsequences $i'_m$, $q'_m-p'_m\to0$.
 We estimate
\[
    TV_K(u_m, Z_j^\delta)
    = \sum_{\ell=0}^{k_m} K(\rho_\ell^m)
\]
with $\rho_\ell^m=q'_{m_\ell}-p'_{m_\ell}$,
\[
    p'_{m_\ell} = u_m(z_{i_m^j+\ell}^m), \quad
    q'_{m_\ell} = u_m(z_{i_m^j+\ell}^m+0), \quad
    i_m + k_m = \bar{i}_m^j, \quad
    \ell=0,\ldots,k_m.
\]
By our assumption of case 2, $\rho_\ell^m\to0$ as $m\to\infty$.
 Moreover, the convergence is uniform, i.e.,
\begin{equation} \label{EUni}
    \lim_{m\to\infty} \max_{0\le\ell\le k_m} \rho_\ell^m = 0.
\end{equation}
Indeed, if otherwise, there were an interval $(p'_{m_\ell},q'_{m_\ell})$ which subsequently converges to an interval with positive length as $m\to\infty$.
 By $L^1$ convergence of $u_m^{-1}$, this convergence is a full convergence which is reduced to Case 1.
 This would contradict $p=q$ so \eqref{EUni} is proved.

By induction (see Proposition~\ref{PInd} below), we observe that
\[
    \sum_{\ell=0}^{k_m} K(\rho_\ell^m)
    \ge K \left(\sum_{\ell=0}^{k_m} \rho_\ell^m \right)
    + C_M \sum_{0\le i<\ell\le k_m} \rho_i^m \rho_\ell^m, \quad
    M \ge 2\sum_{\ell=0}^{k_m} \rho_\ell^m.
\]
Since
\[
    \sum_{\ell=0}^{k_m} \rho_\ell^m
    = q_m - p_m =: s_m
\]
and $s_m\to\rho_j$, we see that
\[
    \sum_{0\le i<\ell\le k_m} \rho_i^m \rho_\ell^m
    = \frac12 \left( s_m^2 - \sum_{\ell=0}^{k_m} (\rho_\ell^m)^2 \right) \to \frac12 \rho_j^2
    \quad\text{as}\quad m \to \infty
\]
since
\[
    \sum_{\ell=0}^{k_m} (\rho_\ell^m)^2
    \le \max_{1\le\ell\le k_m} \rho_\ell^m \cdot s_m \to 0
\]
by \eqref{EUni}.
 Thus
\[
    TV_K(u,Z_j^\delta) \ge K(\rho_j) + \frac{C_M \rho_j^2}{2}.
\]
Again, this contradicts \eqref{EATVK} by taking $\varepsilon$ smaller than $C_M\rho_j^2/2$.
 We thus conclude that the Case 2 does not occur.

We now conclude the existence of a sequence $\left\{I_{\lambda_{i_m^j}^m}(\rho_{i_m^j}^m)\right\}$ converging to $I_{\lambda_j}(\rho_j)$ as $m\to\infty$.

It remains to prove the last statement.
 If such a sequence exists, by $L^1$-convergence of $u_m^{-1}$ the limit must be one of $I_{\lambda_j}(\rho_j)$ if the size does not converge to zero.
 By \eqref{EUni}, the proof of Theorem~\ref{TCJ} is now complete.
\end{proof}

\end{document}
